\section{Алгоритм с изменяемой длительностью измерительного окна, равной периоду инжектируемого сигнала}\label{sec:N_var}

Данный алгоритм характеризуется тем, что в серии из $K$ изменений периода ${T_{z}}$ инжектируемого сигнала (${T^{(k)}_{z}}$, ${k \in \left[1;K\right]}$):
\begin{itemize}
    \item Длительность ${T_h}$ измерительного окна равна периоду ${T_z}$ инжектируемого сигнала (${T^{(k)}_h = T^{(k)}_z}$, ${T^{(k)}_h - var}$).
    \item Период дискретизации ${T_s}$ одинаков для всех значений ${T_z}$ и равен периоду ${T_d}$ выполнения операций в системе (${T^{(k)}_s = T_d}$, ${T^{(k)}_s = const}$).
    \item Количество точек $N$ дискретной последовательности, по которым вычисляются коэффициенты разложения в выражении \eqref{eqn:DFT}, изменяется пропорционально ${T_z}$ (${N^{(k)} - var}$).
\end{itemize}

Период инжектируемого сигнала
\begin{equation*}
    T^{(k)}_z = N^{(k)} \mathord{\cdot} T^{(k)}_s = N^{(k)} \mathord{\cdot} T_d
\end{equation*}
кратен периоду ${T_d}$ выполнения операций в системе (рисунок~\ref{fig:N_var}).

\begin{figure}[H]
	\centering
	\includegraphics [width=1.00\linewidth] {frequencyResponse/N_var.drawio.png}

	\caption{Размер ${N^{(k)}}$ дискретной последовательности при изменении периода ${T^{(k)}_z}$ инжектируемого сигнала}
	\label{fig:N_var}
\end{figure}

В связи с тем, что ${T^{(k)}_h = T^{(k)}_z}$ составляющей с частотой ${f^{(k)}_z}$ соответствуют коэффициенты ${X^{(k)}\left[l\right]}$, ${Y^{(k)}\left[l\right]}$ всегда при ${l = 1}$, т.е. коэффициент передачи динамического звена на частоте ${f^{(k)}_z}$ определяется выражением \eqref{eqn:complexFreqResp} при ${L^{(k)} = 1}$.

Минимальное количество отсчетов $N$ на измерительном окне ${T_h}$ принято равным 3, поэтому частота ${f_z}$ ограничена сверху значением ${f_d / 3}$.

Таким образом, изложенному алгоритму измерения частотных характеристик с последовательной сменой частоты ${f_z}$ инжектируемого сигнала в серии из $K$ измерений \textbf{характерны следующие особенности}:
\begin{enumerate}
    \item Минимально возможная частота ${f_z}$ напрямую определяется максимальным значением измерительного окна ${T_h}$, которое ограничено сверху только объемом памяти для хранения массивов-откликов ${x\left[n\right]}$, ${y\left[n\right]}$.
    
    \item Максимально возможная частота ${f_z}$ в пределе составляет ${f_d / 3}$.
    
    \item В памяти достаточно хранить ${N^{(k)} = f_d / f^{(k)}_z}$-значений каждого массива-отклика ${x^{(k)}\left[n\right]}$, ${y^{(k)}\left[n\right]}$, причем величина $N$ не является постоянной и уменьшается с ростом частоты ${f_z}$ вплоть до ${N = 3}$ при ${f_z \rightarrow f_d / 3}$.
    
    \item Шаг приращения периода ${T_z}$ инжектируемого сигнала (разрешающая способность): ${\Delta T_z = T_d}$.
    
    \item Потенциальные значения частот ${f_z}$ инжектируемого сигнала неравномерно распределены на интервале ${f \in \left(0; {f_d / 2} \right)}$ возможных значений:
    \begin{itemize}
        \item На полуинтервале ${f \in \left[{f_d / 10}; {f_d / 2} \right)}$ частота ${f_z}$ может принимать значения всего из 8-ми возможных: \\
        ${f_z = \left({f_d / 10}, {f_d / 9}, ..., {f_d / 3} \right)}$.
        \item На полуинтервале ${f \in \left[{f_d / 100}; {f_d / 10} \right)}$ -- одно из 90 возможных: \\
        ${f_z = \left({f_d / 100}, {f_d / 99}, ..., {f_d / 11} \right)}$.
        \item На интервале ${f \in \left(0; {f_d / 100} \right)}$ -- сколь угодно много значений, ограниченное исключительно объемом памяти для хранения массивов-откликов ${x\left[n\right]}$, ${y\left[n\right]}$.
    \end{itemize}
\end{enumerate}